\documentclass[11pt,letterpaper]{article}
\usepackage[T1]{fontenc}
\usepackage[utf8]{inputenc}
\usepackage[scaled=0.95]{helvet}
\renewcommand{\familydefault}{\sfdefault}
\usepackage[top=1.9cm,bottom=1.6cm,left=2.3cm,right=2.3cm]{geometry}
\usepackage{microtype}
\usepackage[hidelinks]{hyperref}
\pagestyle{empty}
\setlength{\parindent}{0pt}
\setlength{\parskip}{8pt}

\begin{document}

{\large\bfseries Grover Aruquipa}\par\vspace{-6pt}
{\small Qu\'ebec, QC, Canada \ \textbullet\ \ grover.aruquipa-aruquipa.1@ulaval.ca \ \textbullet\ \ groveraruquipa.github.io}

\vspace{4pt}
Prof.\ Robert J.\ Wood, Microrobotics Laboratory\\
Harvard John A.\ Paulson School of Engineering and Applied Sciences

Dear Professor Wood,

I am applying for the postdoctoral position in your lab on embodied intelligence in soft and reconfigurable robots. I am finishing my PhD in Mechanical Engineering at Universit\'e Laval with Prof.\ Cl\'ement Gosselin and Prof.\ Patrice Lambert, and I will deposit my thesis at the end of December.

The idea behind my PhD is simple: a robot can do more if part of the intelligence is in its body. I work on parallel robots where the platform itself is reconfigurable, so the same mechanism can move the object and also grasp it, with all the motors fixed at the base. I designed, built and controlled three of these robots, with 4, 7 and 8 actuators. I did each one from the type synthesis and the kinematic model to the working prototype, with real-time control at 500\,Hz, and I like the hands-on part as much as the theory.

What attracts me to your lab is that you take the same idea much further, with soft, folded and reconfigurable bodies that are designed together with how they move and sense. I would like to bring my tools to that work: type synthesis to choose the right architecture for a task, and kinematic and singularity analysis with screw theory to understand what a reconfigurable body can and cannot do. I also did clean-room fabrication projects during my master's at FEMTO-ST in micro-robotics, so I am used to working at small scales.

I like working with people. I supervise two graduate students day to day, and I led the Laval team that won first place in the ASME Student Mechanisms and Robotics Design Competition in 2025. After the postdoc I want to start my own lab, and learning from your group how to build new kinds of robot bodies is exactly the step I need.

I am available from early 2027. My CV, two representative papers and three references are attached, and you can see my robots on my website, groveraruquipa.github.io.

Thank you for your time,

Grover Aruquipa

\end{document}
